% SPDX-License-Identifier: MIT
% From h0ffmann/closed-source format/tex/template.tex at 6b3aae5; one change: the author line prints the ORCID iD.
% Paper Fabric article template for pandoc. Compiles with XeLaTeX (the local and CI build) and
% with pdfLaTeX, which is the engine arXiv runs on submitted sources (just arxiv). Modelled on the layout of Sandy Maguire's
% book templates (isovector/thinking-with-types, format/tex/template.tex): one hand-written
% template, metadata from the paper's metadata.md, and the Git revision stamped on every page.
% The CSL and table blocks follow pandoc's default LaTeX template (pandoc -D latex).
\documentclass[10pt,a4paper,]{article}

\usepackage{iftex}
\ifPDFTeX
\usepackage[T1]{fontenc}
\usepackage[utf8]{inputenc}
\usepackage{lmodern}
\usepackage{textcomp}
\usepackage{amssymb}
\else
\usepackage{fontspec}
\usepackage{unicode-math}
\defaultfontfeatures{Ligatures=TeX}
\fi
\usepackage[top=14mm,bottom=17mm,left=18mm,right=18mm]{geometry}
\usepackage{microtype}
\usepackage{xcolor}
\definecolor{fabric}{HTML}{C2410C}
\usepackage{graphicx}
\usepackage{calc}
\usepackage{fancyhdr}
\usepackage{titling}
\usepackage{enumitem}
\setlist{itemsep=2pt,topsep=4pt}
\usepackage[english]{babel}
\usepackage{csquotes}
\setlength{\parindent}{0pt}
\setlength{\parskip}{6pt plus 2pt minus 1pt}
\setlength{\emergencystretch}{3em}
\providecommand{\tightlist}{\setlength{\itemsep}{0pt}\setlength{\parskip}{0pt}}


\usepackage{longtable,booktabs,array}
\usepackage{etoolbox}
\makeatletter
\patchcmd\longtable{\par}{\if@noskipsec\mbox{}\fi\par}{}{}
\makeatother
\usepackage{footnotehyper}
\makesavenoteenv{longtable}

% citeproc citations and the reference list, as in pandoc 3.7's common.latex (pandoc >= 3.1.8
% writes \citeproc and \bibitem). The \small in the list is ours.
\NewDocumentCommand\citeproctext{}{}
\NewDocumentCommand\citeproc{mm}{%
  \begingroup\def\citeproctext{#2}\cite{#1}\endgroup}
\makeatletter
 \let\@cite@ofmt\@firstofone
 \def\@biblabel#1{}
 \def\@cite#1#2{{#1\if@tempswa , #2\fi}}
\makeatother
\newlength{\cslhangindent}
\setlength{\cslhangindent}{1.5em}
\newlength{\csllabelwidth}
\setlength{\csllabelwidth}{3em}
\newenvironment{CSLReferences}[2] % #1 hanging-indent, #2 entry-spacing
 {\small\begin{list}{}{%
  \setlength{\itemindent}{0pt}
  \setlength{\leftmargin}{0pt}
  \setlength{\parsep}{0pt}
  \ifodd #1
   \setlength{\leftmargin}{\cslhangindent}
   \setlength{\itemindent}{-1\cslhangindent}
  \fi
  \setlength{\itemsep}{#2\baselineskip}}}
 {\end{list}}
\newcommand{\CSLBlock}[1]{\hfill\break\parbox[t]{\linewidth}{\strut\ignorespaces#1\strut}}
\newcommand{\CSLLeftMargin}[1]{\parbox[t]{\csllabelwidth}{\strut#1\strut}}
\newcommand{\CSLRightInline}[1]{\parbox[t]{\linewidth - \csllabelwidth}{\strut#1\strut}}
\newcommand{\CSLIndent}[1]{\hspace{\cslhangindent}#1}

\setcounter{secnumdepth}{3}

\AtBeginEnvironment{longtable}{\small}
\newfontfamily\markfont{DejaVuSans.ttf}
\catcode`\⚠=\active
\protected\def⚠{{\markfont\char"26A0\relax}}
\makeatletter
\@ifpackageloaded{subcaption}{}{\usepackage{subcaption}}
\@ifpackageloaded{caption}{}{\usepackage{caption}}
\captionsetup[subfigure]{margin=0.5em}
\AtBeginDocument{%
\renewcommand*\figurename{Figure}
\renewcommand*\tablename{Table}
}
\AtBeginDocument{%
\renewcommand*\listfigurename{List of Figures}
\renewcommand*\listtablename{List of Tables}
}
\newcounter{pandoccrossref@subfigures@footnote@counter}
\newenvironment{pandoccrossrefsubfigures}{%
\setcounter{pandoccrossref@subfigures@footnote@counter}{0}
\begin{figure}\centering%
\gdef\global@pandoccrossref@subfigures@footnotes{}%
\DeclareRobustCommand{\footnote}[1]{\footnotemark%
\stepcounter{pandoccrossref@subfigures@footnote@counter}%
\ifx\global@pandoccrossref@subfigures@footnotes\empty%
\gdef\global@pandoccrossref@subfigures@footnotes{{##1}}%
\else%
\g@addto@macro\global@pandoccrossref@subfigures@footnotes{, {##1}}%
\fi}}%
{\end{figure}%
\addtocounter{footnote}{-\value{pandoccrossref@subfigures@footnote@counter}}
\@for\f:=\global@pandoccrossref@subfigures@footnotes\do{\stepcounter{footnote}\footnotetext{\f}}%
\gdef\global@pandoccrossref@subfigures@footnotes{}}
\@ifpackageloaded{float}{}{\usepackage{float}}
\floatstyle{ruled}
\@ifundefined{c@chapter}{\newfloat{codelisting}{h}{lop}}{\newfloat{codelisting}{h}{lop}[chapter]}
\floatname{codelisting}{Listing}
\newcommand*\listoflistings{\listof{codelisting}{List of Listings}}
\makeatother

\usepackage{xurl}
\usepackage{hyperref}
\hypersetup{
  pdftitle={Wave models on GPUs},
  pdfauthor={Matheus Hoffmann; Pedro Veras Guimarães; Pedro Cruz},
  pdfkeywords={WAVEWATCH III, bit-for-bit
reproducibility, Kokkos, Rust, GPU porting, coding agents},
  colorlinks=true, linkcolor=fabric, citecolor=fabric, urlcolor=fabric,
  pdfcreator={Paper Fabric (pandoc + XeLaTeX)}}
\urlstyle{same}

% The build writes the short Git hash and commit count here (justfile, build/.current-revision).
\newcommand{\revision}{\IfFileExists{.current-revision}{\input{.current-revision}}{unversioned}}

\pagestyle{fancy}
\fancyhf{}
\renewcommand{\headrulewidth}{0pt}
\fancyfoot[L]{\footnotesize\color{gray}Wave models on
GPUs \textperiodcentered{} v0 (preprint, not peer reviewed)}
\fancyfoot[R]{\footnotesize\color{gray}rev.~\revision{} \textperiodcentered{} \thepage}

\begin{document}

\begin{center}
{\LARGE\bfseries Wave models on GPUs\par}
\vspace{4pt}
{\large C++/Kokkos and Rust ports of WAVEWATCH III that must give the
Fortran's bits\par}
\vspace{12pt}
{\large Matheus
Hoffmann}\,\href{https://orcid.org/0009-0009-1056-7661}{\footnotesize\color{fabric}ORCID 0009-0009-1056-7661}\\
{\small\itshape Escola Politécnica, Universidade Federal do Rio de
Janeiro, Brazil}\par\vspace{4pt}
{\large Pedro Veras
Guimarães}\,\href{https://orcid.org/0000-0002-1686-4722}{\footnotesize\color{fabric}ORCID 0000-0002-1686-4722}\\
{\small\itshape Co-advisor of the first author. LabECO, Departamento de
Engenharia Mecânica, Universidade Federal de Santa Catarina,
Brazil}\par\vspace{4pt}
{\large Pedro
Cruz}\,\href{https://orcid.org/0000-0002-2019-133X}{\footnotesize\color{fabric}ORCID 0000-0002-2019-133X}\\
{\small\itshape Advisor of the first author. Grupo de Teleinformática e
Automação (GTA), DEL, Escola Politécnica, Universidade Federal do Rio de
Janeiro, Brazil}\par\vspace{4pt}
\vspace{6pt}
{\small 11 November 2026 \textperiodcentered{} v0 (preprint, not peer
reviewed) \textperiodcentered{} rev.~\revision\par}
\end{center}

\begin{quote}
\small\textbf{Abstract.} We report two ports of the WAVEWATCH III (WW3)
spectral wave model. Coding agents wrote both ports, and both ports have
one acceptance rule: the port must give the same bits as the Fortran. A
tolerance is a diagnostic, never the verdict. \emph{ww3-gpu} keeps WW3,
and the Fortran calls C++ written with Kokkos, behind one run-time flag
for each routine. Its reference is the Fortran path of the same
executable. \emph{chaceon} replaces WW3 with Rust, compiled for the CPU,
for SPIR-V (rust-gpu) and for PTX (Rust-CUDA). No Fortran runs in its
process, and its reference is a WW3 that is compiled separately. ww3-gpu
ports the discrete interaction approximation (DIA) of the four-wave
transfer. On the CPU, it is byte-identical on all 100 serial and 25
OpenMP entries of the WW3 regression matrix that select it. The DIA
takes 48.4 \% of single-thread self time on a case of 9 520 points.
chaceon runs \texttt{ww3\_grid}, \texttt{ww3\_strt}, \texttt{ww3\_prep}
and \texttt{ww3\_shel} without Fortran, and 96 output files are
byte-identical to the reference. Neither port has run the whole model on
a real GPU, and this note claims no speed-up. We compare how each
project defines and checks its evidence.

\textbf{Keywords:} WAVEWATCH III; bit-for-bit
reproducibility; Kokkos; Rust; GPU porting; coding agents.
\end{quote}

\setcounter{section}{-1}

\section{Context: an ocean model ported by a language
model}\label{sec:context}

Koldunov et al. (\citeproc{ref-koldunov2026}{Koldunov et al. 2026})
report the port of a whole ocean model by a large language model. The
model is FESOM2, an ocean and sea-ice model, and the authors are at the
Alfred Wegener Institute. Claude Code, with Claude Opus 4.7, translated
it from Fortran to C and then from C to C++/Kokkos. The people set the
strategy, reviewed the plans and the diffs, and found physics and
validation errors. The agent read, translated, built the test harnesses,
debugged, profiled and wrote the documentation.

The C stage is the centre of their method. The clean C code is the
reference, and each Kokkos kernel keeps its C original as a ``twin''.
The acceptance rule changes at each step. The C model is compared with
the Fortran by statistics: over five years, the RMS difference of the
mean is 0.006 °C in sea-surface temperature. Only the Kokkos Serial
model is byte-identical, and only to the C model. The authors write that
bit-for-bit agreement between Fortran and C ``across compilers and
arithmetic orderings is neither attainable nor necessary''. On the GPU,
they accept differences up to a recorded ceiling for each field.

The successor of WW3 is also in progress. NCEP Office Note 525
(\citeproc{ref-on525}{Tolman 2025}) chose a bottom-up rewrite, with a
C++ core first and a Rust core later ⚠. The same report demotes Fortran
to a solver-only language ⚠. The WAVEWATCH IV repository
(\citeproc{ref-ww4}{NOAA Environmental Modeling Center 2026}) is
pre-alpha: on 2026-11-11 its \texttt{develop} branch had 40 commits,
from 2025-11-17 to 2026-10-16 (\texttt{875f49a}). Its 97 source files
are C++ (66 \texttt{.cpp}, 31 \texttt{.h}), and no physics kernel of WW3
is in it yet.

Thus the core kernels of the wave model stay in Fortran for some years
⚠. This is the first reason why the first author will continue to
maintain ww3-gpu and chaceon. They make the Fortran kernels of WW3 run
on GPUs now, without a change to their bits.

This note starts from that result and asks a narrower question. Can an
agent port a spectral wave model, and give the bits of the Fortran
itself? The two ports below have no C stage, and their reference is the
Fortran. They also use two practices that FESOM2 reports: a strictly
literal translation, and a \texttt{file:line} with the literal value for
each statement about the Fortran.

\section{Why bit for bit}\label{sec:why}

WW3 integrates the action balance equation

\[\frac{\partial N}{\partial t} + \nabla\cdot(\dot{\mathbf{c}}\,N) = \frac{S}{\sigma},\]

where the source term \(S\) is the sum of wind input, dissipation,
nonlinear transfer and bottom terms. A port that changes the order of a
floating-point sum changes \(N\). A forecast system that was validated
against observations must then be validated again. Thus both projects
keep the arithmetic of the Fortran. They do not reorder a reduction,
they do not use a fused multiply-add (\texttt{-ffp-contract=off} on both
sides), and they do not widen a tolerance. A difference is a defect of
the port until two runs of the Fortran itself give different bytes in
the same file.

\section{Two shapes for one rule}\label{sec:shapes}

\begin{longtable}[]{@{}
  >{\raggedright\arraybackslash}p{(\linewidth - 4\tabcolsep) * \real{0.2200}}
  >{\raggedright\arraybackslash}p{(\linewidth - 4\tabcolsep) * \real{0.4000}}
  >{\raggedright\arraybackslash}p{(\linewidth - 4\tabcolsep) * \real{0.3800}}@{}}
\caption{\label{tbl:shapes}The two ports. Sources: chaceon
\texttt{docs/WW3\_GPU\_PARITY\_202611.md} (2026-11-09) and
\texttt{README.md} at \texttt{0c098fe} (\citeproc{ref-chaceon}{Hoffmann
2026a}); ww3-gpu \texttt{kokkos/PORT\_STATUS.md} at \texttt{58167d7}
(\citeproc{ref-ww3gpu}{Hoffmann 2026c}).}\tabularnewline
\toprule\noalign{}
\begin{minipage}[b]{\linewidth}\raggedright
\end{minipage} & \begin{minipage}[b]{\linewidth}\raggedright
ww3-gpu (C++/Kokkos)
\end{minipage} & \begin{minipage}[b]{\linewidth}\raggedright
chaceon (Rust)
\end{minipage} \\
\midrule\noalign{}
\endfirsthead
\toprule\noalign{}
\begin{minipage}[b]{\linewidth}\raggedright
\end{minipage} & \begin{minipage}[b]{\linewidth}\raggedright
ww3-gpu (C++/Kokkos)
\end{minipage} & \begin{minipage}[b]{\linewidth}\raggedright
chaceon (Rust)
\end{minipage} \\
\midrule\noalign{}
\endhead
\bottomrule\noalign{}
\endlastfoot
Shape & Inside WW3: the Fortran calls C++ through a \texttt{bind(C)}
shim, one flag for each routine & Replaces WW3: no Fortran in a chaceon
process \\
Reference & NOAA-EMC/WW3 \texttt{03adda91}, through a fork that adds 5
687 lines and removes none & MetOcean WW3 7.14, \texttt{92e2cf66} \\
CPU parity against & the Fortran path of the same executable & a
reference that is compiled separately \\
Device targets & Kokkos 5.2.2, CUDA backend (\texttt{sm\_89}); OpenMP on
the CPU & rust-gpu (SPIR-V, Vulkan) and Rust-CUDA (PTX) \\
Spectrum resident on the device & yes, on the OpenMP backend & yes, in
SPIR-V on CPU Vulkan (lavapipe) \\
Size of the port & 62 257 lines of C++, 355 files (2026-10-14) & 29 229
lines of Rust port, 5 936 lines of tools \\
Formal methods & SMT queries, exhaustive sweeps & one Lean 4 proof about
\texttt{f32} division \\
Licence & MIT, LGPL-3.0 kernels & MPL-2.0 \\
\end{longtable}

\begin{longtable}[]{@{}
  >{\raggedright\arraybackslash}p{(\linewidth - 4\tabcolsep) * \real{0.3053}}
  >{\raggedright\arraybackslash}p{(\linewidth - 4\tabcolsep) * \real{0.3474}}
  >{\raggedright\arraybackslash}p{(\linewidth - 4\tabcolsep) * \real{0.3474}}@{}}
\caption{\label{tbl:repos}The two repositories. Source: the GitHub REST
API (\texttt{repos/h0ffmann/\textless{}repo\textgreater{}} for
\texttt{created\_at} and \texttt{pushed\_at}, \texttt{commits?sha=main}
for the commits), read on 2026-11-11 at 03:35 UTC. The first commit of
ww3-gpu is older than the repository because it was made before the
first push.}\tabularnewline
\toprule\noalign{}
\begin{minipage}[b]{\linewidth}\raggedright
Metric
\end{minipage} & \begin{minipage}[b]{\linewidth}\raggedright
ww3-gpu
\end{minipage} & \begin{minipage}[b]{\linewidth}\raggedright
chaceon
\end{minipage} \\
\midrule\noalign{}
\endfirsthead
\toprule\noalign{}
\begin{minipage}[b]{\linewidth}\raggedright
Metric
\end{minipage} & \begin{minipage}[b]{\linewidth}\raggedright
ww3-gpu
\end{minipage} & \begin{minipage}[b]{\linewidth}\raggedright
chaceon
\end{minipage} \\
\midrule\noalign{}
\endhead
\bottomrule\noalign{}
\endlastfoot
Repository created & 2026-09-12 01:31 UTC & 2026-10-23 03:58 UTC \\
First commit on \texttt{main} & 2026-09-11 13:47 UTC, \texttt{b43a6d3} &
2026-10-23 03:58 UTC, \texttt{3bad31a} \\
Last commit on \texttt{main} & 2026-11-11 03:16 UTC, \texttt{efc0211} &
2026-11-09 08:53 UTC, \texttt{0c098fe} \\
Last push, any branch & 2026-11-11 03:31 UTC & 2026-11-09 08:53 UTC \\
Commits on \texttt{main} & 800 & 110 \\
\end{longtable}

The Kokkos route keeps every WW3 program and replaces one routine at a
time. Thus one regtest can isolate the effect of one routine. The Rust
route must rebuild the programs, the file formats and the time loop
before it can compare a whole run. After that, no Fortran remains that
can share state with the port.

The CPU results of the whole ww3-gpu runtime, with the resident
spectrum, and their limits are in its NOAA progress record
(\citeproc{ref-noaarecord}{Hoffmann 2026e}). The DIA port is also
published alone as ww3-dia-gpu-v2 (\citeproc{ref-diagpuv2}{Hoffmann
2026d}). That repository compares the port with an earlier CUDA kernel
of the same routine.

\section{Evidence ladders}\label{sec:ladders}

Each project gives names to its steps. A step in one ledger is not the
same evidence as a step in the other ledger.

\begin{longtable}[]{@{}
  >{\raggedright\arraybackslash}p{(\linewidth - 4\tabcolsep) * \real{0.2400}}
  >{\raggedright\arraybackslash}p{(\linewidth - 4\tabcolsep) * \real{0.3800}}
  >{\raggedright\arraybackslash}p{(\linewidth - 4\tabcolsep) * \real{0.3800}}@{}}
\caption{\label{tbl:ladders}Evidence steps. ww3-gpu counts come from
\texttt{scripts/port\_score.py\ -\/-write} (profile
\texttt{bench\_small}, 120 × 80 points, 12 h, one thread, gprof,
2026-10-14); chaceon counts come from
\texttt{just\ port-summary}.}\tabularnewline
\toprule\noalign{}
\begin{minipage}[b]{\linewidth}\raggedright
Step
\end{minipage} & \begin{minipage}[b]{\linewidth}\raggedright
ww3-gpu
\end{minipage} & \begin{minipage}[b]{\linewidth}\raggedright
chaceon
\end{minipage} \\
\midrule\noalign{}
\endfirsthead
\toprule\noalign{}
\begin{minipage}[b]{\linewidth}\raggedright
Step
\end{minipage} & \begin{minipage}[b]{\linewidth}\raggedright
ww3-gpu
\end{minipage} & \begin{minipage}[b]{\linewidth}\raggedright
chaceon
\end{minipage} \\
\midrule\noalign{}
\endhead
\bottomrule\noalign{}
\endlastfoot
Code exists & \emph{translated}: 316 routines, 26.2 \% of Fortran lines
& \emph{translated}: Rust exists, not fully reached \\
A run calls it & \emph{ported}: 72 routines, 99.3 \% of profiled run
time & \emph{routine}: every captured call is bit-identical on the host
and in SPIR-V \\
Whole runs identical, CPU & \emph{regtested}: 14 routines, 48.6 \% of
run time & \emph{host}: whole runs are \texttt{cmp}-identical; 10 080
lines \\
Identical through a GPU API & none & \emph{spirv}: the same files from
SPIR-V on CPU Vulkan; 5 189 lines \\
Identical on a real GPU & \emph{gpu}: none & \emph{cuda}: none \\
\end{longtable}

\section{What each project found}\label{sec:found}

\textbf{ww3-gpu: the reference is the optimised code.} Byte identity for
the DIA needed one fact that the source text does not give. gfortran
expands \texttt{x**11} as one chain of multiplications at \texttt{-O0}
and as a different chain at \texttt{-O3}. The two chains disagree on 643
427 of 1 000 000 single-precision inputs. Before this was found, the
research log explained the remaining differences as round-off that WW3
amplifies. That explanation was wrong. Thus the reference of the port is
the code that the WW3 build produces, read as GIMPLE or disassembly, not
the Fortran text. The port has no unit tests of its own. Only the WW3
regression tests, compared by \texttt{test.comp}, judge it.

\textbf{chaceon: a run with no Fortran in it.} The Rust programs
\texttt{ww3\_grid}, \texttt{ww3\_strt}, \texttt{ww3\_prep} and
\texttt{ww3\_shel} run in sequence from the namelists. They write
\texttt{mod\_def.ww3}, the forcing, the gridded and point outputs, the
track file and both restarts. All 96 files are byte-identical to the
reference. The boundaries give 214 identical comparisons over 34 cases,
and the track post-processors give 74.
\texttt{ww3\_shel\ -\/-threads\ N} and an MPI build give the same bytes.
The files stay identical when \texttt{W3SRCE}, \texttt{W3XYP3} and
\texttt{W3KTP3} are compiled to SPIR-V and run on a CPU Vulkan driver.
Enumeration cannot reach the \(2^{64}\) pairs of \texttt{f32} operands.
There, a Lean proof states the division fact that the port uses. The
proof says nothing about whether Rust equals Fortran, because Fortran
has no usable formal semantics.

\textbf{Why ww3-gpu has no proof yet, and what comes next.} In this
first phase, ww3-gpu adopted no proof assistant
(\citeproc{ref-adr0001}{Hoffmann 2026f}). There are four reasons:

\begin{enumerate}
\def\labelenumi{\arabic{enumi}.}
\tightlist
\item
  No proof assistant reads Fortran or C++. A proof is about a model that
  a person writes. It says nothing about the gfortran and nvcc binaries
  unless somebody shows that the model agrees with them, and that
  agreement is the open question.
\item
  Bit equality is a property of two binaries, built with given flags. It
  is not a property of two source texts.
\item
  For straight-line IEEE-754 code, equality on every input is decidable.
  SMT queries in the \texttt{QF\_FP} theory and sweeps over all
  \(2^{32}\) \texttt{float32} inputs answer it in seconds to minutes. An
  interactive proof takes months for each routine.
\item
  Agents must write each proof, and a person must review it. This is a
  second artefact for each routine.
\end{enumerate}

Without a proof, the DIA port has the strongest evidence short of one.
It is byte-identical on all 125 entries of the WW3 regression matrix
that select it (100 serial, 25 OpenMP). Section 1 of the kernel was also
swept over all \(2^{32}\) \texttt{float32} inputs, with 0 differences
(\citeproc{ref-adr0001}{Hoffmann 2026f}). Thus the DIA port is validated
by the tests of the model that it replaces, on every entry that calls
it. This validation is not a formal proof, and the note does not call it
one.

A first formal layer is already next to the ports. The Lean 4 package
\texttt{ocean-waves} (\citeproc{ref-oceanwaves}{Hoffmann 2026b}) states
the discretised action balance equation and proves these properties,
with no \texttt{sorry} and no \texttt{axiom}:

\begin{itemize}
\tightlist
\item
  \(\sigma^2(k)\) increases strictly with \(k\);
\item
  \(c/2 \le c_g \le c\);
\item
  the depth CFL condition keeps every Courant number in \([0, 1]\);
\item
  upwind propagation and a DIA-shaped quadruplet exchange conserve
  action;
\item
  the semi-implicit source step keeps the action positive for every time
  step.
\end{itemize}

The package uses exact \texttt{Rat} arithmetic on a reduced model: one
space dimension, a fixed depth and a Padé model of \(\tanh\). Thus it
also says nothing about bits. Agda is the next step. Agda has no 32-bit
float type, so it cannot prove a fact about bits. Its role is a total,
executable reference semantics of each source term. Each table index
will have a \texttt{Fin} type, so every read is in bounds for every
grid, not only for one grid. The sweeps, the solver queries and the WW3
regression tests continue to decide the bits.

\section{Challenges: the regression tests are
large}\label{sec:challenges}

WW3's regression tests are the only judge of the Kokkos port, and they
are large. Table Table~\ref{tbl:regtests} gives their size at the fork
commit \texttt{80644ba9}, counted on 2026-10-18.

\begin{longtable}[]{@{}
  >{\raggedright\arraybackslash}p{(\linewidth - 4\tabcolsep) * \real{0.7158}}
  >{\raggedright\arraybackslash}p{(\linewidth - 4\tabcolsep) * \real{0.2000}}
  >{\raggedright\arraybackslash}p{(\linewidth - 4\tabcolsep) * \real{0.0842}}@{}}
\caption{\label{tbl:regtests}Size of the WW3 regression matrix. Source:
ww3-gpu \texttt{regtests/README.md} at \texttt{58167d7}, counts from
\texttt{just\ ww3-gpu-matrix-list}. The count of all entries is the
number of lines that \texttt{matrix.base} writes, and a machine section
can disable some of them.}\tabularnewline
\toprule\noalign{}
\begin{minipage}[b]{\linewidth}\raggedright
Set of \texttt{matrix.base} entries
\end{minipage} & \begin{minipage}[b]{\linewidth}\raggedright
Entries
\end{minipage} & \begin{minipage}[b]{\linewidth}\raggedright
Share
\end{minipage} \\
\midrule\noalign{}
\endfirsthead
\toprule\noalign{}
\begin{minipage}[b]{\linewidth}\raggedright
Set of \texttt{matrix.base} entries
\end{minipage} & \begin{minipage}[b]{\linewidth}\raggedright
Entries
\end{minipage} & \begin{minipage}[b]{\linewidth}\raggedright
Share
\end{minipage} \\
\midrule\noalign{}
\endhead
\bottomrule\noalign{}
\endlastfoot
All entries & 758, in 60 tests & 100 \% \\
Entries with the DIA (\texttt{NL1}) & 310 & 41 \% \\
Entries with ST4 physics (\texttt{ST4}) & 136 & 18 \% \\
Entries that the resident preset can select & 7 (all serial) & 0.9 \% \\
Selected regtests of ww3-gpu & 5, and the \texttt{ww3\_ts1} pair & 0.7
\% \\
\end{longtable}

One run of the 310 \texttt{NL1} entries takes days. The selected set
takes one to three hours on one workstation (Intel i9-14900, 32 hardware
threads); this time is an estimate, not a measurement ⚠. Thus a full
matrix run cannot follow each commit. ww3-gpu runs the selected set as a
guardrail, and it postpones the full matrix to a validation phase. Until
then, a static argument supports bit parity: the \texttt{-O3} GIMPLE or
the disassembly of both compilers.

\section{What is not claimed}\label{sec:limits}

\begin{enumerate}
\def\labelenumi{\arabic{enumi}.}
\tightlist
\item
  Neither project has run the whole model on a real GPU. The PTX of
  chaceon has never run. ww3-gpu has run one DIA benchmark kernel on an
  RTX 4090 (\citeproc{ref-diagpuv2}{Hoffmann 2026d}).
\item
  This note claims no speed-up. The rules of ww3-gpu mark its recorded
  timings (\citeproc{ref-noaarecord}{Hoffmann 2026e}) ``regtest
  pending'', and chaceon has no timings.
\item
  On a GPU, bit identity also needs the \texttt{expf} and \texttt{logf}
  of the device to return the bits of glibc. ww3-gpu sends its kernels
  to device copies of the glibc 2.44 functions. These copies equal the
  host functions on every input that was checked. No device run has
  compared them yet. Until one does, a fixed tolerance judges a GPU run.
\item
  The two references are different lines of WW3 (NOAA-EMC and MetOcean).
  Thus the comparison of the two ports is by routine name, not by source
  text. Both ledgers contain 85 routines.
\end{enumerate}

\section{Use of AI}\label{sec:ai}

Claude coding agents (Anthropic) produced the C++/Kokkos and Rust code,
the shims, the comparison scripts, the investigations and drafts of this
text. They worked under the direction of the first author, in the two
repositories. Pedro Cruz (GTA, UFRJ) is the advisor and Pedro Veras
Guimarães (LabECO, UFSC) is the co-advisor of the first author.

Most of the commits came from Claude Opus. Claude Sonnet also ported a
large amount of code without gross errors, in the experience of the
first author ⚠. One example is the split of \texttt{W3SDS4} into parts
that run one sea point per vector lane (ww3-gpu commit \texttt{9cef132},
841 lines added). Its output files were byte-identical to the Fortran of
the same executable on the \texttt{smoke} and \texttt{full} cases.

Each repository publishes the rules, skills and hooks of its agents in
its \texttt{AGENTS.md} file and its \texttt{.claude/} directory. Every
number in this note comes from a file of a repository at the commit that
the note names. Nothing was run for this note. Each claim was checked on
2026-11-11 against the source and the commit that it names; a ⚠ marks a
claim that was not checked against a primary source. The physics of WW3
is used without change; WAVEWATCH III is a registered trademark of NOAA.

\textbf{Disclaimer.} Today, both repositories are set up for Claude
Code. The authors will continue to work so that any coding agent can use
ww3-gpu and chaceon. This includes agents that run local models with
Ollama. This work is not complete, and this note does not claim results
for other agents or for local models.

\section*{References}\label{references}
\addcontentsline{toc}{section}{References}

\protect\phantomsection\label{refs}
\begin{CSLReferences}{1}{0}
\bibitem[\citeproctext]{ref-chaceon}
Hoffmann, Matheus. 2026a. {``{chaceon}: A {Rust} Port of {WAVEWATCH
III}.''} \url{https://github.com/h0ffmann/chaceon}.

\bibitem[\citeproctext]{ref-oceanwaves}
---------. 2026b. {``Ocean Waves: The Physics of a Spectral Wave Model
in {Lean} 4.''}
\url{https://github.com/h0ffmann/proof/tree/b863c15/ocean-waves}.

\bibitem[\citeproctext]{ref-ww3gpu}
---------. 2026c. {``{WW3 GPU Lab}: Agentic Coding and Agentic Research
on a Bit-for-Bit {GPU} Port of {WAVEWATCH III}.''}
\url{https://doi.org/10.5281/zenodo.23221351}.

\bibitem[\citeproctext]{ref-diagpuv2}
---------. 2026d. {``{ww3-dia-gpu-v2}: The Whole {W3SNL1}, Inside
{WAVEWATCH III}, Bit for Bit.''}
\url{https://github.com/h0ffmann/ww3-dia-gpu-v2/tree/32aa2ef}.

\bibitem[\citeproctext]{ref-noaarecord}
---------. 2026e. {``A Bit-Identical {Kokkos} Runtime for {WAVEWATCH
III}: {CPU} Results and Five Speed-Ups.''} ww3-gpu,
\texttt{docs/NOAA\_PROGRESS\_RECORD\_202610.md}, commit 58167d7.
\url{https://github.com/h0ffmann/ww3-gpu/blob/58167d7/docs/NOAA_PROGRESS_RECORD_202610.md}.

\bibitem[\citeproctext]{ref-adr0001}
---------. 2026f. {``{ADR-0001}: No Proof Language for the {GPU} Port;
{SMT} Queries and Exhaustive Sweeps Instead.''} ww3-gpu,
\texttt{docs/ADRs/ADR-0001-proof-language.md}, commit 58167d7.
\url{https://github.com/h0ffmann/ww3-gpu/blob/58167d7/docs/ADRs/ADR-0001-proof-language.md}.

\bibitem[\citeproctext]{ref-koldunov2026}
Koldunov, Nikolay V., Suvarchal K. Cheedela, Sergey Danilov, Dmitry
Sidorenko, Sebastian Beyer, and Thomas Jung. 2026. {``An Ocean Model
Ported by a Large Language Model: Experience and Lessons from {FESOM2}
({Fortran} to {C} to {C++/Kokkos}).''} arXiv:2606.11356
{[}physics.ao-ph{]}. \url{https://doi.org/10.48550/arXiv.2606.11356}.

\bibitem[\citeproctext]{ref-ww4}
NOAA Environmental Modeling Center. 2026. {``{WAVEWATCH IV}.''}
\url{https://github.com/NOAA-EMC/WW4}.

\bibitem[\citeproctext]{ref-on525}
Tolman, Hendrik L. 2025. {``The {WAVEWATCH III} Software Modernization
Project: Phase {I} Report.''} Office Note 525. College Park, MD:
NOAA/NWS/NCEP. \url{https://doi.org/10.25923/h7j3-1h25}.

\end{CSLReferences}

\end{document}
